\subsection{El teorema de pureza de Van der Waerden para el conjunto de
ramificación de un morfismo birracional}
\label{subsection:IV.21.12}
\label{IV.21.12}

\begin{env}[21.12.1]
\label{IV.21.12.1}
Sean $X$ y $U$ dos preesquemas, y sea $f:U\to X$ un morfismo cuasicompacto y
cuasiseparado, de modo que $f_*(\mathcal O_U)$ es una $\mathcal O_X$-álgebra
cuasicoherente (1.7.4). Se llama \emph{envolvente afín} del $X$-preesquema $U$ al
$X$-preesquema afín sobre $X$
\[
  U^0=\operatorname{Aff}(U/X)=\operatorname{Spec}(f_*(\mathcal O_U))
  =\operatorname{Spec}(\mathscr A(U))\qquad(\mathbf{II}, 1.1.1).
\]
Si $f^0:U^0\to X$ es el morfismo estructural, se tiene, por definición,
\[
  \mathscr A(U^0)=f^0_*(\mathcal O_{U^0})=f_*(\mathcal O_U)=\mathscr A(U),
\]
y al isomorfismo identidad de $f_*(\mathcal O_U)$ le corresponde, por
$(\mathbf{II}, 1.2.7)$, un $X$-morfismo canónico
\[
\label{IV.21.12.1.1}
  i_U:U\to U^0.
  \tag{21.12.1.1}
\]

Para que $i_U$ sea un \emph{isomorfismo}, es evidentemente necesario y suficiente que el
morfismo $f:U\to X$ sea \emph{afín}. Para todo $X$-preesquema $V$ \emph{afín sobre
$X$}, la aplicación $u\mapsto u\circ i_U$ es una \emph{biyección}
\[
  \operatorname{Hom}_X(U^0,V)\xrightarrow{\sim}\operatorname{Hom}_X(U,V)
\]
functorial en $V$: esto resulta de la existencia de las biyecciones canónicas
\[
  \operatorname{Hom}_X(U,V)\xrightarrow{\sim}
  \operatorname{Hom}_{\mathcal O_X}(\mathscr A(V),\mathscr A(U))
\]
y
$\operatorname{Hom}_X(U^0,V)\xrightarrow{\sim}
\operatorname{Hom}_{\mathcal O_X}(\mathscr A(V),\mathscr A(U^0))$
$(\mathbf{II}, 1.2.7)$.

Se puede decir, pues, que $U^0$ \emph{representa} el funtor covariante
$V\mapsto\operatorname{Hom}_X(U,V)$ en la categoría de los $X$-preesquemas afines sobre
$X$ $(\mathbf 0_{\mathrm{III}}, 8.1.11)$. Se deduce de ello
$(\mathbf 0_{\mathrm{III}}, 8.1.7)$ que, para $X$ fijo,
$U\mapsto\operatorname{Aff}(U/X)$ es un funtor covariante de la categoría de los
$X$-preesquemas cuasicompactos y cuasiseparados sobre $X$ en la categoría de los
$X$-preesquemas afines sobre $X$; más precisamente, si $U_1$, $U_2$ son dos
$X$-preesquemas cuasicompactos y cuasiseparados sobre $X$, a todo $X$-morfismo
$g:U_1\to U_2$ le corresponde el único $X$-morfismo $g^0:U_1^0\to U_2^0$ que hace
conmutativo el diagrama
\[
  \xymatrix{
    U_1 \ar[r]^g \ar[d]_{i_{U_1}} & U_2 \ar[d]^{i_{U_2}} \\
    U_1^0 \ar[r]_{g^0} & U_2^0.
  }
\]
\end{env}

\oldpage[IV-4]{305}
Más generalmente, consideremos un diagrama conmutativo
\[
  \xymatrix{
    U \ar[r]^v \ar[d]_f & U' \ar[d]^{f'} \\
    X \ar[r]_u & X'
  }
\]
donde los morfismos $f$, $f'$ son cuasicompactos y cuasiseparados y el morfismo $u$ es
\emph{afín}. Si se pone $h=u\circ f$, se tiene
$h_*(\mathcal O_U)=u_*(f_*(\mathcal O_U))=u_*(f^0_*(\mathcal O_{U^0}))$ y
$u\circ f^0$ es un morfismo afín; se tiene, por consiguiente,
$U^0=\operatorname{Aff}(U/X')$ (relativamente al morfismo $h$), de donde resulta un único
$X'$-morfismo $v^0:U^0\to U'^0$ que hace conmutativo el diagrama
\[
  \xymatrix{
    U \ar[r]^v \ar[d]_{i_U} & U' \ar[d]^{i_{U'}} \\
    U^0 \ar[r]_{v^0} & U'^0.
  }
\]

\begin{proposition}[21.12.2]
\label{IV.21.12.2}
Sean $f:U\to X$ un morfismo cuasicompacto y cuasiseparado y
$h:X'\to X$ un morfismo plano; se pone $U'=U\times_X X'$,
$f'=f_{\{X'\}}:U'\to X'$. Entonces se tiene un $X'$-isomorfismo canónico
\[
\label{IV.21.12.2.1}
  \operatorname{Aff}(U'/X')\xrightarrow{\sim}
  \operatorname{Aff}(U/X)\times_X X'.
  \tag{21.12.2.1}
\]
\end{proposition}

\begin{proof}
En efecto, se tiene
$\operatorname{Aff}(U'/X')=\operatorname{Spec}(f'_*(\mathcal O_{U'}))$, y
$\operatorname{Aff}(U/X)\times_X X'=\operatorname{Spec}(h^*(f_*(\mathcal O_U)))$
$(\mathbf{II}, 1.5.1)$; el isomorfismo del enunciado procede del isomorfismo canónico
$h^*(f_*(\mathcal O_U))\xrightarrow{\sim}f'_*(\mathcal O_{U'})$ (2.3.1).
\end{proof}

\begin{corollary}[21.12.3]
\label{IV.21.12.3}
Para todo morfismo cuasicompacto y cuasiseparado $f:U\to X$ y todo $x\in X$, se tiene,
salvo isomorfismo canónico,
\[
  U^0\times_X\operatorname{Spec}(\mathcal O_{X,x})
  =(U\times_X\operatorname{Spec}(\mathcal O_{X,x}))^0.
\]
\end{corollary}

Resulta asimismo de (21.12.2) que se tiene, salvo isomorfismo canónico, para todo
abierto $V$ de $X$,
\[
\label{IV.21.12.4}
  (f^{-1}(V))^0=(f^0)^{-1}(V).
  \tag{21.12.4}
\]

\begin{env}[21.12.5]
\label{IV.21.12.5}
Consideraremos, en particular, el caso en que $f:U\to X$ es una \emph{inmersión
abierta}, identificándose así $U$ con un preesquema inducido sobre un abierto de $X$.
Como el
morfismo $f^{-1}(U)\to U$ es la identidad, resulta de (21.12.4) que el morfismo
$(f^0)^{-1}(U)\to U$, restricción de $f^0$, es un isomorfismo, siendo entonces
$i_U:U\to U^0$ una \emph{inmersión abierta} que permite identificar $U$ con un
preesquema inducido sobre un abierto de $U^0$.

Más generalmente, para un abierto $V$ de $X$, la restricción $(f^0)^{-1}(V)\to V$ de
$f^0$ es un isomorfismo de $(f^0)^{-1}(V)$ sobre el preesquema inducido sobre el abierto
$V\cap U$ si y solo si la inmersión abierta $U\cap V\to V$ es un morfismo
\emph{afín}. Es claro $(\mathbf{II}, 1.2.1)$
\oldpage[IV-4]{306}
que la unión de estos abiertos $V$ es el \emph{mayor} de ellos, $U_1$, que
contiene a $U$; en virtud de lo precedente, $U_1$ es también el mayor abierto que no
interseca el conjunto
\[
  \operatorname{Daf}(U/X)=f^0(U^0)-U
\]
(«defecto de afinidad» del abierto $U$ relativamente a $X$, que es vacío si y solo
si $U$ es afín sobre $X$); en otras palabras, el conjunto cerrado $Z=X-U_1$ es la
\emph{adherencia} del conjunto $\operatorname{Daf}(U/X)$.

Obsérvese que, para todo morfismo \emph{plano} $h:X'\to X$, si se pone
$U'=h^{-1}(U)$, se tiene
\[
\label{IV.21.12.5.1}
  \operatorname{Daf}(U'/X')=h^{-1}(\operatorname{Daf}(U/X))
  \tag{21.12.5.1}
\]
como resulta inmediatamente de (21.12.2.1) y de $(\mathbf I, 3.4.8)$. En particular,
para todo abierto $V$ de $X$, se tiene
\[
\label{IV.21.12.5.2}
  \operatorname{Daf}((U\cap V)/V)=\operatorname{Daf}(U/X)\cap V
  \tag{21.12.5.2}
\]
y, para todo $x\in X$,
\[
\label{IV.21.12.5.3}
  \operatorname{Daf}((U\cap\operatorname{Spec}(\mathcal O_{X,x}))
  /\operatorname{Spec}(\mathcal O_{X,x}))
  =\operatorname{Daf}(U/X)\cap\operatorname{Spec}(\mathcal O_{X,x}).
  \tag{21.12.5.3}
\]
\end{env}

Cuando $X$ sea localmente noetheriano, precisaremos la naturaleza del conjunto
$\operatorname{Daf}(U/X)$, lo cual mostrará, por ejemplo, que cuando $U$ es denso en
todo $X$, $\operatorname{Daf}(U/X)$ no es un conjunto cerrado raro
arbitrario:

\begin{theorem}[21.12.6]
\label{IV.21.12.6}
Sean $X$ un preesquema localmente noetheriano, $U$ un abierto no vacío de $X$ y
$f:U\to X$ la inyección canónica. Entonces:
\begin{enumerate}[label=(\roman*)]
  \item La adherencia $Z=X-U_1$ de $\operatorname{Daf}(U/X)$ es de codimensión
  $\geq2$ en $X$.
  \item Si $T=X-U\supset Z$ es de codimensión $\geq2$, el morfismo
  $f^0:U^0\to X$ es sobreyectivo.
\end{enumerate}
\end{theorem}

\begin{proof}
(i) Demostremos primero que, para todo punto $x\in\operatorname{Daf}(U/X)$, se tiene
necesariamente $\dim(\mathcal O_{X,x})>1$. En efecto, $x$ evidentemente no puede ser un punto
maximal de $X$, pues está contenido en $\overline U-U$; se trata, pues, de ver que no se
puede tener $\dim(\mathcal O_{X,x})=1$. Ahora bien, esta relación implicaría, por (21.12.5.3),
$x\in\operatorname{Daf}(U^{(x)}/\operatorname{Spec}(\mathcal O_{X,x}))$, donde
$U^{(x)}=U\cap\operatorname{Spec}(\mathcal O_{X,x})$. Pero los únicos abiertos de
$\operatorname{Spec}(\mathcal O_{X,x})$ son el propio $\operatorname{Spec}(\mathcal O_{X,x})$
y los subconjuntos del conjunto (finito) de los puntos maximales de
$\operatorname{Spec}(\mathcal O_{X,x})$. Ahora bien, el conjunto abierto reducido a un punto maximal
de $\operatorname{Spec}(\mathcal O_{X,x})$ es afín, por definición de los preesquemas; se
concluye de ello que \emph{todos} los conjuntos abiertos de
$\operatorname{Spec}(\mathcal O_{X,x})$ son afines y, por tanto,
$\operatorname{Daf}(U^{(x)}/\operatorname{Spec}(\mathcal O_{X,x}))=\varnothing$
contrariamente a la hipótesis formulada.

Para demostrar (i) hay que probar algo más, a saber, que si $x\in X$ es tal que
$\dim(\mathcal O_{X,x})=1$, existe un entorno abierto $W$ de $x$ en $X$ tal que
$W$ no interseca $\operatorname{Daf}(U/X)$, es decir, tal que la inyección
canónica $f_W:U\cap W\to W$ sea afín. Pero, con las mismas notaciones anteriores,
acabamos de ver que la inyección canónica
$f^{(x)}:U^{(x)}\to\operatorname{Spec}(\mathcal O_{X,x})$ es afín. Evidentemente, se
puede limitar al caso en que $X$ es noetheriano; como el morfismo $f$ es de
presentación finita, la conclusión resulta de (8.10.5, (viii)), aplicado según el
método de (8.1.2, a)).

\oldpage[IV-4]{307}
(ii) Se trata de demostrar que, para todo punto $x\in T$, se tiene $x\in f^0(U^0)$;
demostraremos primero que se puede reducir al caso en que $X=\operatorname{Spec}(A)$, donde
$A$ es un anillo local noetheriano \emph{completo}, y $x$ es el punto cerrado de $X$. Para
ello basta efectuar el cambio de base
$h:X'=\operatorname{Spec}(\widehat{\mathcal O}_{X,x})\to X$; si se pone
$U'=h^{-1}(U)$; $f'=f_{\{X'\}}$ es la inyección canónica $U'\to X'$ y, puesto que el
morfismo $h$ es plano, resulta de (21.12.2) que, si se demuestra que $x$ pertenece a
$f^{\prime 0}(U^{\prime 0})$, se deducirá que $x\in f^0(U^0)$. En virtud de (6.1.1), se ha
efectuado así la reducción buscada.

Sea entonces $X_1$ un subpreesquema cerrado reducido de $X$ cuyo espacio subyacente
sea una componente irreducible de $X$, de dimensión \emph{máxima} entre las que
contienen una componente irreducible de $T$, y pongamos
$U_1=U\cap X_1$, $T_1=T\cap X_1=X_1-U_1$; se tiene entonces
$\operatorname{codim}(T_1,X_1)\geq2$ $(\mathbf 0, 14.2.1)$, y el par $(U_1,X_1)$
satisface las mismas hipótesis que el par $(U,X)$ (pues $X_1$ es el espectro de un
anillo cociente de $A$ y, por tanto, local noetheriano y completo). Se tiene, además, un
diagrama conmutativo
\[
  \xymatrix{
    U_1 \ar[r]^j \ar[d]_{f_1} & U \ar[d]^f \\
    X_1 \ar[r]_i & X
  }
\]
donde $i$ y $j$ son las inyecciones canónicas, siendo $i$, por tanto, un morfismo afín; se
deduce de (21.12.1) la existencia de un morfismo $j^0:U_1^0\to U^0$ que hace conmutativo el
diagrama
\[
  \xymatrix{
    U_1^0 \ar[r]^{j^0} \ar[d]_{f_1^0} & U^0 \ar[d]^{f^0} \\
    X_1 \ar[r]_i & X
  }
\]
y, por consiguiente, para demostrar que $x\in f^0(U^0)$ basta demostrar que
$x\in f_1^0(U_1^0)$. Se puede, pues, sustituir $X$ por $X_1$ y suponer, por consiguiente,
que el anillo $A$ es además \emph{íntegro}. Pero $A$ es cociente de un
anillo noetheriano regular en virtud del teorema de Cohen $(\mathbf 0, 19.8.8$, (i)) y,
como es íntegro, se puede aplicar (5.11.1) con la familia $(x_\alpha)$ reducida al
único punto maximal de $X$; como $\operatorname{codim}(T,X)\geq2$ por hipótesis, se
ve que $f_*(\mathcal O_U)$ es un $\mathcal O_X$-módulo \emph{coherente}, de modo que el
morfismo $f^0:U^0\to X$ es finito; como este morfismo es dominante (pues $U$ es denso en
todas partes), es sobreyectivo $(\mathbf{II}, 6.1.10)$ y, por tanto, $x\in f^0(U^0)$. C.Q.D.
\end{proof}

\begin{corollary}[21.12.7]
\label{IV.21.12.7}
Sean $X$ un preesquema localmente noetheriano, $U$ un subpreesquema inducido sobre un
abierto de $X$ y $j:U\to X$ la inmersión canónica; supongamos que $j$ sea un morfismo
afín. Entonces toda componente irreducible de $T=X-U$ es de codimensión $\leq1$
(y, por tanto, de codimensión $1$ si $U$ es denso en todas partes).
\end{corollary}

\begin{proof}
Supongamos, en efecto, que una de las componentes irreducibles $T_1$ de $T$ sea de
codimensión $\geq2$ en $X$. Sustituyendo, si es preciso, $X$ por un entorno abierto del
punto genérico de $T_1$, se puede suponer que $T$ es irreducible y de codimensión $\geq2$.
Pero entonces la hipótesis
\oldpage[IV-4]{308}
de que $j:U\to X$ es afín implica que $U^0$ se identifica con $U$ y $j^0$ con $j$,
es decir, que $j^0(U^0)=U$; ahora bien, esto contradice la conclusión de (21.12.6), que,
bajo la hipótesis de dimensión formulada, implica que $j^0$ debe ser sobreyectivo.
\end{proof}

\begin{env}[21.12.8]
\label{IV.21.12.8}
Sean $A$ un anillo local noetheriano, $Y=\operatorname{Spec}(A)$, $y$ el único punto
cerrado de $Y$ y $Y'=Y-\{y\}$. Consideremos la condición siguiente:
\begin{itemize}
  \item[\textup{(W)}] \emph{Para todo abierto $U$ de $Y$ contenido en $Y'$, que no
  contenga ninguna componente irreducible de $Y'$ y tal que la inmersión canónica
  $U\to Y'$ sea afín, $U$ es a su vez un abierto afín.}
\end{itemize}
Esta condición se deduce de la siguiente:
\begin{itemize}
  \item[\textup{(W')}] \emph{Para toda parte cerrada $T$ de $Y$, cada una de cuyas componentes
  irreducibles sea de codimensión $1$ en $Y$, y tal que para toda componente
  irreducible $Y_i$ de $Y$ se tenga $Y_i\cap T\ne\{y\}$, el abierto $U=Y-T$ es afín.}
\end{itemize}
En efecto, si se satisface (W') y si el abierto $U$ de $Y$ satisface la hipótesis de (W),
resulta de (21.12.7) que ninguna componente irreducible de $T=Y-U$ puede ser de
codimensión $\geq2$; como, además, $U$ no contiene ninguna de las componentes
irreducibles $Y_i-\{y\}$ de $Y'$, la condición (W') muestra que $U$ es afín.

Obsérvese que la condición (W') se simplifica cuando $Y$ es \emph{irreducible} y es
entonces equivalente a la siguiente:
\begin{itemize}
  \item[\textup{(W'')}] \emph{Para toda parte cerrada irreducible $T$ de $Y$, de
  codimensión $1$ en $Y$, $Y-T$ es un abierto afín.}
\end{itemize}
En efecto, es claro que (W') implica (W'') cuando $Y$ es irreducible, y la
recíproca es cierta considerando las componentes irreducibles de $T$ y observando
que la intersección de un número finito de abiertos afines es un abierto afín
$(\mathbf I, 5.5.6)$.
\end{env}

\begin{env}[21.12.9]
\label{IV.21.12.9}
\emph{Ejemplos}. ---
Si $A$ es un anillo local noetheriano \emph{factorial}, satisface (W') (y
\emph{a fortiori} (W)), puesto que todo ideal primo de altura $1$ es principal
$(\mathbf I, 1.3.6)$. Pero existen anillos locales noetherianos no factoriales que
satisfacen (W''), por ejemplo los de dimensión $\leq1$: en efecto, se observó en la
demostración de (21.12.6, (i)) que \emph{todos} los abiertos de $Y$ son entonces afines.
Por otra parte, utilizando la teoría de la dualidad local (cap.~III,
3\textsuperscript{a} parte), se puede demostrar que todo anillo local noetheriano de dimensión $2$
satisface (W').
\end{env}

El interés de la condición (W) reside en el resultado siguiente:

\begin{proposition}[21.12.10]
\label{IV.21.12.10}
Sean $X$, $Y$ dos preesquemas localmente noetherianos, $g:X\to Y$ un morfismo,
$y$ un punto cerrado de $Y$, $Y'=Y-\{y\}$ y $X'=g^{-1}(Y')$; supongamos que el morfismo
$g':X'\to Y'$, restricción de $g$, sea una inmersión abierta, y que el anillo local
$\mathcal O_{Y,y}$ satisfaga la condición (W) (21.12.8). Entonces, para toda componente
irreducible $Z$ de $X_y=g^{-1}(y)$, o bien $Z$ es de codimensión $\leq1$ en $X$,
o bien su punto genérico está aislado en $Z$. Si $Z$ es localmente de tipo finito sobre
$k(y)$, la segunda alternativa implica que $Z$ está reducido a un solo punto.
\end{proposition}

\begin{proof}
La última afirmación del enunciado resulta de que $Z$ es entonces un preesquema de
Jacobson (10.4.7) y, como el conjunto de los puntos cerrados de $Z$ es entonces denso en
$Z$, el punto genérico de $Z$ solo puede estar aislado en $Z$ si $Z$ está reducido a un
solo punto.

\oldpage[IV-4]{309}
Supongamos que existe en $X_y$ una componente irreducible $Z$ cuyo punto
genérico $z$ no está aislado en $Z$, y tal que $\operatorname{codim}(Z,X)\geq2$.
Como la cuestión es local sobre $X$ y $Y$, puesto que $z$ no está aislado en $Z$, se puede,
sustituyendo $X$ por un entorno abierto de $z$ en $X$, suponer que $X$ e $Y$ son afines,
que $X_y=Z$ es irreducible, no está reducido a un punto y satisface
$\operatorname{codim}(X_y,X)\geq2$. La imagen $g(X-X_y)=U$ es un abierto de $Y'$
isomorfo a $X-X_y$ por hipótesis; demostremos que, sustituyendo si es preciso $X$ por un
entorno abierto de $z$ en $X$, se puede suponer que $U$ no contiene ninguna de las
componentes irreducibles de $Y'$ cuya adherencia contiene a $y$.

En efecto, se puede suponer primero que todos los puntos maximales de $X$ son
generalizaciones de $z$, pues el conjunto de estos puntos es finito; por tanto, el conjunto de los puntos
maximales de $U$ es el conjunto de las imágenes $y_i$ por $g$ de los puntos maximales $x_i$ de
$X$ (ya que, por hipótesis, ningún punto maximal de $X$ puede estar contenido en $X_y$,
pues $\dim(\mathcal O_{X,z})\geq2$). Pongamos $X_i=\overline{\{x_i\}}$. Por
hipótesis, se tiene $z\in X_i$ y, como $z$ no está aislado en $X_y$, existe en
$X_y$ un punto $x\ne z$. Como $X_i\ne Z$, existe en $X_i$ un punto $t_i$ tal
que, si se pone $T_i=\overline{\{t_i\}}$, se tenga $\dim(\mathcal O_{T_i,x})=1$
(10.5.9); esto implica, por hipótesis, $t_i\notin Z$, de modo que $t_i$ no es una
generalización de $z$. Sustituyendo $X$ por $X'=X-\bigcup_iT_i$, se ve que la imagen por
$g$ de $X'-X_y$ no contiene las imágenes $g(t_i)$ y, por tanto, no contiene ninguna de las
componentes irreducibles de $Y'$ cuya adherencia contiene a $y$.

Hecho esto, como $X$ e $Y$ son afines, el morfismo $g:X\to Y$ es afín, y lo
mismo sucede con la restricción $g':X'\to Y'$; puesto que $g'$ es una inmersión abierta,
la inmersión canónica $U\to Y'$ es afín. Pongamos
$Y_1=\operatorname{Spec}(\mathcal O_{Y,y})$, $Y'_1=Y_1-\{y\}=Y'\cap Y_1$,
$U_1=U\cap Y_1$. Lo precedente demuestra que la inmersión canónica $U_1\to Y'_1$ es
afín; por tanto, $U_1$ es un abierto afín de $Y_1$ en virtud de (W), es decir, la
inmersión canónica $U_1\to Y_1$ es afín. Pero, como esta inmersión es de
presentación finita (1.6.2), resulta de (8.10.5, (viii)), aplicado según el método de
(8.1.2, a)), que, restringiendo si es preciso $Y$ a un entorno abierto afín de $y$, se puede
suponer que la inmersión $U\to Y$ es afín. Se concluiría que el abierto $X-X_y$ de
$X$, isomorfo a $U$, sería afín, lo que contradiría (21.12.7); queda así demostrada la
proposición.
\end{proof}

\begin{corollary}[21.12.11]
\label{IV.21.12.11}
Sean $X$, $Y$ dos preesquemas localmente noetherianos, suponiéndose que $X$ es
irreducible; sea $g:X\to Y$ un morfismo localmente de tipo finito, y sea $V$ el mayor
abierto de $X$ tal que la restricción $g|V:V\to Y$ sea un isomorfismo local.
Supongamos que, para todo punto $y\in Y$, el anillo local $\mathcal O_{Y,y}$ satisface la
condición (W) (21.12.8). Entonces toda componente irreducible de $T=X-V$, o bien
es de codimensión $\leq1$, o bien tiene su punto genérico $z$ aislado en
$g^{-1}(g(z))$.
\end{corollary}

\begin{proof}
Pongamos $g(z)=y$. Como la cuestión solo depende de la fibra $g^{-1}(y)$ y del anillo
$\mathcal O_{Y,y}$, mediante un cambio de base se puede limitar al caso en que
$Y=\operatorname{Spec}(\mathcal O_{Y,y})$ y $X$ es afín
$((\mathbf I, 3.6.5)$ y $(\mathbf I, 4.5.5))$; sustituyendo eventualmente $X$ por un
entorno abierto de $z$ en $X$, se puede suponer que $T$ es irreducible; además, se
puede limitar al caso en que $V$ no es vacío. Se puede, pues, suponer que el morfismo $g$ es
separado y que $y$ es cerrado en $Y$; puesto que la restricción $g|V:V\to Y$ es un isomorfismo
local y el abierto no vacío $V$ de $X$ es irreducible, resulta de
$(\mathbf I, 8.2.8)$ que $g|V:V\to Y$ es una inmersión abierta. La restricción
$g|(V\cap X_y):V\cap X_y\to\{y\}=\operatorname{Spec}(k(y))$ es, por tanto, también una
inmersión abierta, lo que demuestra que $V\cap X_y$ es, o bien vacío, o bien está reducido a un punto
$x$ racional
\oldpage[IV-4]{310}
sobre $k(y)$ y, por tanto, cerrado en $X_y$ $(\mathbf I, 6.4.2)$. Sustituyendo nuevamente,
si es preciso, $X$ por un entorno abierto más pequeño de $z$ en $X$, se puede, pues,
limitar al caso en que $V\cap X_y=\varnothing$, es decir, $T\supset X_y$; como, además,
$z\in X_y$ es el punto genérico de $T$, se tiene
$g(T)\subset\overline{\{y\}}=\{y\}$, es decir, $T=X_y$. Se está entonces exactamente
en las condiciones de aplicación de (21.12.10), de donde resulta la conclusión.
\end{proof}

\begin{theorem}[21.12.12]
\label{IV.21.12.12}
\textup{(van der Waerden)}. ---
Sean $X$, $Y$ dos preesquemas localmente noetherianos íntegros, y sea $g:X\to Y$ un
morfismo birracional y localmente de tipo finito. Se supone además que $Y$ es
normal y que, para todo $y\in Y$, el anillo $\mathcal O_{Y,y}$ satisface la condición (W)
de (21.12.8) (condiciones que se cumplen, en particular, cuando $Y$ es localmente factorial
(21.12.9)). Si $V$ es el mayor abierto de $X$ tal que la restricción
$g|V:V\to Y$ sea un isomorfismo local, todas las componentes irreducibles de
$T=X-V$ son de codimensión $1$ en $X$.
\end{theorem}

\begin{proof}
Obsérvese que, puesto que $g$ es birracional, el conjunto abierto $V$ no es vacío; basta,
pues, demostrar que un punto maximal $z$ de $T$ no puede estar aislado en $X_y$,
donde $y=g(z)$. Pero, restringiéndose a un entorno abierto de $z$, se puede
limitar al caso en que $T=X_y$; entonces todas las fibras $g^{-1}(y')$ ($y'\in Y$) serían
vacías o estarían reducidas a un punto, y resultaría del «Main theorem» (8.12.10) que $g$
sería un isomorfismo local, contrariamente a la hipótesis $z\in T$.
\end{proof}

\begin{corollary}[21.12.13]
\label{IV.21.12.13}
Supongamos que se satisfacen las hipótesis de (21.12.12) y, además, que $g$ es
cuasifinito en cada uno de los puntos de $X^{(1)}$ (se recuerda que estos son los puntos en los que
$\dim(\mathcal O_{X,x})=1$); entonces $g$ es un isomorfismo local y, si además $g$ es
separado, $g$ es una inmersión abierta.
\end{corollary}

\begin{proof}
Basta demostrar la primera afirmación, pues la segunda es consecuencia de la
primera y de $(\mathbf I, 8.2.8)$. Todo se reduce a demostrar, con las notaciones de
(21.12.12), que $T=\varnothing$; en caso contrario, un punto genérico $z$ de una
componente irreducible de $T$ pertenecería a $X^{(1)}$ en virtud de (21.12.12) y, por
hipótesis, estaría aislado en $X_y$, donde $y=g(z)$; pero se ha visto en la demostración
de (21.12.12) que esto no es posible.
\end{proof}

\begin{remarks}[21.12.14]
\label{IV.21.12.14}
(i) La conclusión de (21.12.12) se aplica cuando $Y$ es regular e íntegro y $X$ es
íntegro, pues, en virtud del teorema de Auslander--Buchsbaum (21.11.1), $Y$ es entonces
localmente factorial. En cambio, se conocen ejemplos en los que $X$ e $Y$ son esquemas
algebraicos normales de dimensión $3$, sobre un cuerpo algebraicamente cerrado de
característica arbitraria, y en los que la conclusión de (21.12.12) ya no es válida.

(ii) El conjunto $T$ de (21.12.12) es también el conjunto de los puntos en los que $g$ está
\emph{ramificado}. En efecto, si $g$ no está ramificado en un punto $x\in X$, tampoco lo está
en un entorno abierto afín $U$ de $x$ (17.3.7), de modo que $g|U$ es un morfismo
separado, cuasifinito y birracional (17.4.3). Puesto que $Y$ es normal, se concluye del «Main
theorem» $(\mathbf{III}, 4.4.9)$ que $g$ es un isomorfismo local en el punto $x$ y, por tanto,
$x\in X-T$. Recíprocamente, $g$ evidentemente no está ramificado en ningún punto en el que es un
isomorfismo local. Esto justifica el título dado a esta sección.

\oldpage[IV-4]{311}
(iii) Sin suponer que $Y$ sea localmente factorial, pero suponiendo, en cambio, que
$X$ es intersección completa relativamente a $Y$ (cap.~V), veremos en el cap.~V
que se tiene un resultado más preciso que (21.12.12), expresando $T$ como ciclo
1-codimensional de una sección de un $\mathcal O_X$-módulo invertible, canónicamente
asociada a $g$. Esto se aplicará, en particular, cuando $X$ e $Y$ sean ambos regulares,
o cuando $g$ sea un $S$-morfismo, siendo $X$ e $Y$ preesquemas lisos sobre $S$.

(iv) Cabe preguntarse si (21.12.10) admite una recíproca; en otras palabras, dado un
anillo local noetheriano $A$, si se pone $Y=\operatorname{Spec}(A)$, designando $y$
el punto cerrado de $Y$, y si, para todo preesquema localmente noetheriano $X$ y
todo morfismo $g:X\to Y$ tal que $g':X'\to Y'$ sea una inmersión abierta, toda
componente irreducible de $X_y$ es de codimensión $\leq1$ en $X$ o está reducida a
un punto, ¿es entonces cierto que $A$ satisface la condición (W) de (21.12.8)?

(v) Sean $Y$ un preesquema íntegro \emph{regular}, $X$ un preesquema localmente
noetheriano \emph{normal} y $g:X\to Y$ un morfismo localmente de tipo finito; supongamos
además que, si $\eta$ es el punto genérico de $Y$, la fibra $X_\eta$ sea \emph{étale}
sobre $k(\eta)$, y sea $T'$ el complementario del mayor abierto $V'$ de $X$ tal que
$g|V':V'\to Y$ sea étale; ¿es cierto entonces que todas las componentes irreducibles de
$T'$ sean de codimensión $1$? Así ocurre cuando se supone además que $g$
es \emph{localmente cuasifinito} («teorema de pureza» de Zariski--Nagata). Se puede
demostrar que la conjetura precedente sería consecuencia de la siguiente (y sería
incluso equivalente a ella si se satisficiera la conjetura de (iv)): para un anillo local regular
$A$ contenido en un anillo local íntegro e íntegramente cerrado $B$, tal que $B$ sea un
$A$-módulo finito, el abierto $U$ de los puntos de $\operatorname{Spec}(B)$ en los que
$\operatorname{Spec}(B)$ no está ramificado sobre $\operatorname{Spec}(A)$ (o, lo que
equivale por (18.10.1), es étale sobre $\operatorname{Spec}(A)$) es afín.
\end{remarks}

\begin{proposition}[21.12.15]
\label{IV.21.12.15}
Sean $S$ un preesquema, $g:X\to S$ un morfismo plano y localmente de presentación
finita, cuyas fibras $X_s=g^{-1}(s)$ son geométricamente irreducibles (4.5.2), y
$h:Y\to S$ un morfismo liso; se denotan por $Y_s=h^{-1}(s)$ las fibras de este morfismo.
Sea $f:X\to Y$ un $S$-morfismo propio tal que, para todo punto maximal $\eta$ de
$S$, el morfismo $f_\eta:X_\eta\to Y_\eta$ sea un isomorfismo. Entonces $f$ es un
isomorfismo.
\end{proposition}

\begin{proof}
Puesto que $h$ es plano, todo punto maximal de $Y$ está sobre un punto maximal de $S$
(2.3.4), de modo que la unión de los $Y_\eta$, cuando $\eta$ recorre el conjunto de los puntos
maximales de $S$, es densa en $Y$; como los $f_\eta$ son isomorfismos, $f$
es dominante y, por tanto, \emph{sobreyectivo}, puesto que es propio. Basta, pues, demostrar
que $f$ es una \emph{inmersión abierta}. Teniendo en cuenta (17.9.5), basta probar
que, para todo $s\in S$, $f_s:X_s\to Y_s$ es una inmersión abierta.

Como todo $s\in S$ es una especialización de un punto maximal $\eta$\footnote{Corrección de fuente: el francés impreso y el inglés mantenido dicen que $s$ es una generalización de $\eta$. La construcción siguiente toma $s\in\overline{\{\eta\}}$, por lo que $s$ es una especialización de $\eta$, o equivalentemente $\eta$ una generalización de $s$.}, mediante el
cambio de base $\operatorname{Spec}(\mathcal O_{S',s})\to S$, donde $S'$ es el
subpreesquema reducido de $S$ cuyo espacio subyacente es $\overline{\{\eta\}}$,
se puede reducir desde ahora al caso en que $S$ es un esquema local e íntegro de punto cerrado $s$:
en efecto, las fibras en los puntos $\eta$ y $s$ no cambian y las propiedades de los
morfismos $g$, $h$ y $f$ se conservan por cambio de base. Además, la cuestión
es local sobre $Y$; sustituyendo $Y$ por un entorno abierto afín $U$ en $Y$ de un punto
de $Y_s$, y $X$ por $f^{-1}(U)$, que es cuasicompacto puesto que $f$ es propio, se ve
que se puede suponer, además, que $X$ e $Y$ son de presentación finita sobre $S$.

Existen entonces un subanillo local noetheriano $A_0$ de
$A=\mathcal O_{S,s}$ tal que $A_0\to A$ sea un homomorfismo local, dos morfismos
de tipo finito $g_0:X_0\to S_0=\operatorname{Spec}(A_0)$,
$h_0:Y_0\to S_0$ y un $S_0$-morfismo
\oldpage[IV-4]{312}
$f_0:X_0\to Y_0$ tales que $g$, $h$ y $f$ se deducen de $g_0$, $h_0$, $f_0$ por el
cambio de base $S\to S_0$ ((8.9.1) y (5.13.3)); se puede, además, suponer que $g_0$ es
plano (11.2.7), $h_0$ liso (17.7.9) y $f_0$ propio (8.10.5, (xii)). Por otra parte, la
proyección sobre $S_0$ del punto genérico $\eta$ de $S$ es el punto genérico $\eta_0$ de
$S_0$, y resulta de (2.7.1, (viii)) que el morfismo
$(f_0)_{\eta_0}:(X_0)_{\eta_0}\to(Y_0)_{\eta_0}$ es un isomorfismo; por último, el punto
cerrado $s$ de $S$ está sobre el punto cerrado $s_0$ de $S_0$, de modo que la fibra
$(X_0)_{s_0}$ de $g_0$ es geométricamente irreducible (4.5.6).

Se ve, pues, que se puede limitar al caso en que $S$ es el espectro de un anillo local íntegro
noetheriano, de punto cerrado $s$, debilitar la hipótesis sobre las fibras suponiendo solamente
que $X_s$ y $X_\eta$ son geométricamente irreducibles, y demostrar que, bajo estas
hipótesis, $f_s$ es una inmersión abierta. Existe entonces un anillo de valoración discreta
$A'$ y un morfismo $S'=\operatorname{Spec}(A')\to S$ que transforma el punto cerrado
$a$ de $S'$ en $s$ y el punto genérico $b$ de $S'$ en $\eta$ $(\mathbf{II}, 7.1.9)$;
sean $g':X'\to S'$, $h':Y'\to S'$, $f':X'\to Y'$ los morfismos deducidos de $g$,
$h$, $f$ por el cambio de base $S'\to S$, que satisfacen de nuevo las hipótesis
satisfechas por $g$, $h$ y $f$ en el enunciado (21.12.15); si se demuestra que
$f'_a:X'_a\to Y'_a$ es una inmersión abierta, resultará de (2.7.1, (x)) que también
lo es $f_s:X_s\to Y_s$. Se puede, pues, limitar finalmente al caso en que $S$ es
el \emph{espectro de un anillo de valoración discreta}.

Entonces, puesto que $h:Y\to S$ es liso, $Y$ es regular (17.5.8), y como la cuestión
es local sobre $Y$, se puede limitar al caso en que $Y$ es íntegro. Como $g:X\to S$ es
plano, los puntos maximales de $X$ están contenidos en $X_\eta$ (2.3.4) y, puesto que
$f_\eta$ es un isomorfismo, $X$ es irreducible y el anillo local en su punto
genérico es un cuerpo; además, por planitud (3.3.2), se ve que $X$ no tiene ciclos
primos inmersos, de modo que $X$ es reducido (3.2.1), por tanto íntegro, y $f$ es un
morfismo birracional y separado. Para demostrar que $f$ es una inmersión abierta, se puede
aplicar, pues, el criterio (21.12.13); para ver que $f$ es cuasifinito en los puntos de
$X^{(1)}$, se observa que la afirmación es evidente en aquellos de estos puntos que
pertenecen a $X_\eta$; el único punto de $X^{(1)}$ que pertenece a $X_s$ es el punto
genérico $x$ de $X_s$, en virtud de (6.1.1). Ahora bien, resulta de (2.4.6) y de (14.2.2)
que los morfismos $g$ y $h$ son equidimensionales; como $X_\eta$ e $Y_\eta$ son
isomorfos, las componentes irreducibles de $X_s$ y de $Y_s$ tienen todas la misma
dimensión. Pero, por hipótesis, $X_s$ es irreducible y se ha visto que $f$ es sobreyectivo,
por lo que también lo es $f_s:X_s\to Y_s$, lo cual implica que $Y_s$ es también
irreducible; si $y$ es el punto genérico de $Y_s$, se tiene, pues, $f(x)=y$, y la relación
$\dim(X_s)=\dim(Y_s)$ implica entonces, en virtud de (5.6.6), que
$\dim(f_s^{-1}(y))=0$, es decir, $f_s$ (y, por tanto, también $f$) es efectivamente cuasifinito
en el punto $x$, lo que termina la demostración.
\end{proof}

\begin{corollary}[21.12.16]
\label{IV.21.12.16}
Sean $g:X\to S$ un morfismo propio, plano y de presentación finita, $h:Y\to S$ un
morfismo liso, propio y de presentación finita, y $f:X\to Y$ un $S$-morfismo. Se
supone que las fibras $X_s=g^{-1}(s)$ de $g$ son geométricamente irreducibles. Sea
$U$ el conjunto de los $s\in S$ tales que $f_s:X_s\to Y_s$ sea un isomorfismo. Entonces
$U$ es abierto y cerrado en $S$, y el morfismo $g^{-1}(U)\to h^{-1}(U)$,
restricción de $f$, es un isomorfismo.
\end{corollary}

\begin{proof}
La última afirmación resultará de la primera y de (17.9.5). Ya se sabe (9.6.1,
(xi)) que $U$ es localmente constructible en $S$. En virtud de (1.10.1), basta
demostrar que $U$ es estable por \emph{especialización} y por \emph{generalización}. Para
demostrar estas afirmaciones se
\oldpage[IV-4]{313}
puede, mediante un cambio de base $\operatorname{Spec}(\mathcal O_{S,s})\to S$,
reducir al caso en que $S$ es un esquema local de punto cerrado $s$ y punto genérico
$\eta$, y hay que demostrar que, para que $f_s$ sea un isomorfismo, es necesario y suficiente
que $f_\eta$ lo sea. Ahora bien, la suficiencia de esta condición resulta de (21.12.15). Para
demostrar que es necesaria, se razona como en la demostración de (21.12.15)
(observando, con las mismas notaciones, que la proyección del punto cerrado de $S$ es el
punto cerrado de $S_0$) para reducir al caso en que, además, $S$ es \emph{noetheriano};
pero entonces la conclusión resulta de $(\mathbf{III}, 4.6.7$, (ii)).
\end{proof}

\begin{remarks}[21.12.17]
\label{IV.21.12.17}
(i) La conclusión de la proposición (21.12.15) ya no es válida si se elimina
la hipótesis de que las fibras $X_s$ son irreducibles. Tomemos, en efecto,
$S=\operatorname{Spec}(A)$, donde $A$ es un anillo de valoración discreta, y
$Y=\mathbb P^1_S$, que es propio y liso sobre $S$ (17.3.9). Denotemos nuevamente por $s$ y
$\eta$ el punto cerrado y el punto genérico de $S$; sea $z$ un punto cerrado de $Y_s$,
por ejemplo un punto racional sobre $k=k(s)$, y sea $X$ el $Y$-preesquema obtenido por
explosión del punto $z$. Como el anillo de polinomios $A[T]$ es regular
$(\mathbf 0, 17.3.7)$ y de dimensión $2$, se ve, como en la demostración de
(15.1.1.6), que, si $f:X\to Y$ es el morfismo estructural, $f^{-1}(z)$ es isomorfo
a $\mathbb P^1_k$ y, por otra parte, el complementario de $f^{-1}(z)$ en $X_s$ es
isomorfo a $Y_s-\{z\}$; por tanto, $Z_1=f^{-1}(z)$ y $Z_2$, adherencia en $X_s$ del
complementario de $Z_1$, son las dos componentes irreducibles de $X_s$. Por
otra parte, $f$ es evidentemente propio y $g=h\circ f$ es plano, pues $X$ es íntegro
$(\mathbf{II}, 8.1.4)$ y, para todo abierto afín $U$ de $X$, el homomorfismo
$A\to\Gamma(U,\mathcal O_X)$ es inyectivo $(\mathbf I, 1.2.7)$, de modo que
$\Gamma(U,\mathcal O_X)$ es un $A$-módulo sin torsión y, por tanto, plano, puesto que $A$
es un anillo de valoración discreta $(\mathbf 0_{\mathrm I}, 6.3.4)$. Es claro que
$f_\eta:X_\eta\to Y_\eta$ es un isomorfismo, aunque
$f_s:X_s\to Y_s$ no lo sea.

(ii) La conclusión de (21.12.15) también falla si, entre las hipótesis, se
suprime la hipótesis de que $f$ es \emph{propio}. Para verlo, basta definir $S$ e
$Y$ como en (i), pero sustituir $X$ por el complementario $X'$ de $Z_2$ en el
preesquema $X$ definido en (i), y $f$ por la restricción $f':X'\to Y$ de $f$; el morfismo
$g':X'\to S$, restricción de $g$, sigue siendo plano, y esta vez $X'_s$ es
geométricamente irreducible; además, $X'_s$ es isomorfo al complementario de un
punto cerrado en $\mathbb P^1_k$, por lo que es un esquema afín isomorfo a $\mathbb A^1_k$;
como su imagen en $Y_s$ está reducida al punto cerrado $z$, $f'$ no es propio
$(\mathbf{III}, 4.4.2)$ y $f'_s$ no es un isomorfismo, aunque $f'_\eta$ sí lo sea.

(iii) Es posible que el enunciado de la proposición (21.12.15) siga siendo válido cuando se
sustituya en él la palabra «isomorfismo» por «morfismo étale» (cf. (21.12.14, (v))). En
tal caso, lo mismo seguirá siendo cierto para (21.12.16).
\end{remarks}
